\statement{B.106}{Proved}{A free double quotient preserves the actual selected fixed-source question in the discrete and full profinite categories}
\label{stmt:B106}
Use the actual based surface $Z$, involution $F$ and fixed source
$S$ of B.87--B.105. Explicitly, put $U=\mu _5\setminus\{1\}$ and
$\mathscr B=\mathbb C\setminus(\{0\}\cup\mu _5)$. The surface $Z$
consists of $(v,w)\in(\mathbb C^*)^2$ with
\[
 v,w\ne1,\quad v\ne w,\quad q=vw\in\mathscr B,
\]
and with the following four equations excluded for each $u\in U$:
\[
 v=u,\qquad w=u,\qquad q+(u-1)v-u=0,\qquad q+(u-1)w-u=0.
\]
Its base point is $P=(399/200,1/200)$, its literal involution is
\[
 F(v,w)=\left(\frac{v-q}{v-1},\frac{w-q}{w-1}\right),
 \qquad S=Z\cap\{v+w=2\}.
\]
The free coordinate swap $\sigma(v,w)=(w,v)$ commutes with $F$.
The analytic quotient $\mathscr Q=Z/\langle\sigma\rangle$ is exactly
the following open surface in $\mathscr B\times\mathbb P^1_t$:
\begin{equation}\label{eq:B106-surface}
 t\ne1,-1,\qquad qt^2\ne1,\qquad
 t\ne T_u(q),-T_u(q)\quad(u\in U),\qquad
 T_u(q)=\frac{q+u^2-2u}{u^2+(1-2u)q}.
\end{equation}
The graph equations are projective, including their poles. The
coordinates and descended involution are
\begin{equation}\label{eq:B106-chart}
 s=v+w,\qquad t=\frac{s-2}{s-2q},\qquad
 s=\frac{2(qt-1)}{t-1},\qquad F_{\mathscr Q}(q,t)=(q,-t).
\end{equation}
The based quotient point is $\bar P=(399/40000,0)$.

Put $\alpha_u=u(2-u)$ and $\beta_u=u^2/(2u-1)$. The two fixed
sources in $\mathscr Q$ are
\[
 D_0=(\mathscr B\setminus\{\alpha_u:u\in U\})\times\{0\},
 \qquad
 D_\infty=(\mathscr B\setminus\{\beta_u:u\in U\})\times\{\infty\}.
\]
Each is a ten-punctured affine line. The actual map $S\to D_0$
is a connected unramified double cover. In particular its source
group is $F_{19}$, whereas $\pi_1(D_0)=F_{10}$.

Let $\Pi_Z=\pi_1(Z,P)$ and $\Pi_Q=\pi_1(\mathscr Q,\bar P)$,
and let $I_Z,I_Q$ be the respective images of these selected source
groups. Then, for the literal based actions,
\begin{equation}\label{eq:B106-fixed-equivalence}
 \begin{aligned}
  \Pi_Z^F=I_Z&\quad\Longleftrightarrow\quad\Pi_Q^F=I_Q,\\
  (\widehat{\Pi_Z})^F=\overline{I_Z}
    &\quad\Longleftrightarrow\quad
       (\widehat{\Pi_Q})^F=\overline{I_Q}.
 \end{aligned}
\end{equation}
Here completions are full profinite completions, and every bar is
the closure in the indicated ambient completion. No source injection
or global fixed-source equality is assumed or concluded.

\begin{proof}
\emph{B.106.1. The exact quotient.}
At fixed $q$, the two maps on $v$ are
$f(v)=(v-q)/(v-1)$ and $g(v)=q/v$. Direct substitution gives
$f^2=g^2=1$ and $fg=gf$. The swap has no fixed point on $Z$.
The invariant coordinates $s,q$ recover the unordered roots of
$X^2-sX+q$, and their discriminant is removed. Thus the actual
quotient map is an unramified covering of degree two.
The descended map in the $s$ coordinate is
\[
 s\longmapsto\frac{4q-(1+q)s}{q+1-s}.
\]
Its two fixed sections are $s=2,2q$, distinct on $\mathscr B$.
The change \eqref{eq:B106-chart} sends it to $t\mapsto-t$.
The missing affine point $s=\infty$ gives $t=1$; the deleted
$s=q+1$, where $v=1$ or $w=1$, gives $t=-1$.
Moreover
\[
 s^2-4q=\frac{4(q-1)(qt^2-1)}{(t-1)^2}.
\]
The pair $v=u$ or $w=u$ has $s=u+q/u$ and hence $t=T_u(q)$.
Its $F$ image is $t=-T_u(q)$. Conversely these are all the old
deleted factors. This proves both directions of
\eqref{eq:B106-surface}, including the allowed points at $t=\infty$.

For comparison with the previous reflection chart, its coordinates
satisfy
\[
 t=-b/a,\qquad a^2=\frac{q-1}{qt^2-1},\qquad b=-ta.
\]
At infinity use $b^2=(q-1)/q\ne0$; no branch is introduced there.
The generic quotient fiber has twelve punctures on a projective
line, and its actual double cover has twenty-two, agreeing with
the ranks $11$ and $21$ of the previous fiber groups.

\emph{B.106.2. The complete collision formulas.}
Write $d_u=u^2+(1-2u)q$. Direct polynomial calculations give
\begin{align*}
 d_u(T_u-1)&=2u(q-1),&
 d_u(T_u+1)&=2(1-u)(q-u),\\
 d_u^2(qT_u^2-1)&=(q-1)(q-u^2)^2,\\
 d_ud_r(T_u-T_r)&=2(u-r)(q-1)(q-ur),\\
 d_ud_r(T_u+T_r)&=-2(q+ur-r-u)\bigl(q(r+u-1)-ur\bigr).
\end{align*}
For distinct $u,r$, the same-sign collisions and the incidences
with the branch curve or the retained punctures lie over already
deleted fifth roots. All additional graph collisions are opposite-sign
ones, at
\[
 q=u+r-ur\quad\hbox{or}\quad q=ur/(u+r-1).
\]
The denominator $u+r-1$ is nonzero: if two unit complex numbers
sum to $1$, they are $e^{\pm i\pi/3}$, which are not fifth roots.
Taking $u=r$ gives exactly the deleted points $\alpha_u$ on $D_0$
and $\beta_u$ on $D_\infty$.

We have $|\alpha_u|=|2-u|>1$ and
$\beta_u=1/\overline{\alpha_u}$. Equality of two alpha values gives
$(u-r)(2-u-r)=0$, and two unit numbers sum to $2$ only when both
are $1$. Thus the four alpha values are distinct and outside the
unit circle, and the four beta values are distinct and inside it.
They are nonzero and avoid all six old punctures. This proves the
ten-puncture counts. Over $D_0$ the cover is obtained by adjoining
$\sqrt{1-q}$; monodromy around the omitted $q=1$ is nontrivial.
It is therefore connected, and its index-two source group has
rank $2(10-1)+1=19$. Its inverse image is exactly $S$.
The other fixed source lifts to the fixed locus of $F\sigma$,
which is not substituted for $S$.

\emph{B.106.3. Based group and full-topology transfer.}
Identify $\Pi_Z$ with the index-two subgroup of $\Pi_Q$ supplied
by the cover, and let $\varepsilon:\Pi_Q\to C_2$ be its monodromy.
The pullback source cover proves
\[
 I_Q\cap\Pi_Z=I_Z,\qquad \varepsilon(I_Q)=C_2.
\]
The first assertion does not require a source injection: a source
loop has image in $\Pi_Z$ exactly when its cover monodromy is
trivial, equivalently when it lifts to a closed loop in $S$.
The second follows from the loop about $1$. Both source images
are fixed pointwise. An element fixed in $\Pi_Q$ either belongs
to $\Pi_Z$ or becomes such an element after multiplication by a
fixed element of $I_Q$ with the same nontrivial monodromy.
This proves the discrete equivalence.

A finite-index subgroup carries its full profinite topology:
the core in the ambient group of any finite-index subgroup of it
again has finite index. Consequently $\widehat{\Pi_Z}$ is the
open kernel of $\widehat\varepsilon$ in $\widehat{\Pi_Q}$.
The same finite-index argument applies to the source cover.
The images of the compact completed source groups are their
actual closures, and hence
\[
 \overline{I_Q}\cap\widehat{\Pi_Z}=\overline{I_Z},
 \qquad \widehat\varepsilon(\overline{I_Q})=C_2.
\]
The identical two-coset argument proves the full profinite
equivalence. In particular no completion of an infinite-index
source image has been silently identified with its closure.
The retained program \texttt{code/swap\_quotient21.py} verifies
the rational identities and actual base point; its two initial
failed attempts and the successful exact output are all retained.
\end{proof}

\statement{B.107}{Proved}{The selected fixed-source equality holds across every critical fiber in the actual inner region}
\label{stmt:B107}
In B.106 put $\mathscr B_< =\{q:0<|q|<1\}$, and let
$\mathscr Q_<,Z_<$ be the actual inverse images of this base.
No critical fiber in this region is removed. Let $S_<^+$ be the
component of $S\cap Z_<$ containing $P$, and let
$D_< =\mathscr B_<\times\{0\}\subset\mathscr Q_<$.
For either $(X,L)=(\mathscr Q_<,D_<)$ or $(Z_<,S_<^+)$,
the based source map is injective with its full profinite topology,
and
\begin{equation}\label{eq:B107-local-fixed}
 \pi_1(X)^F=\pi_1(L)\cong\mathbb Z,
 \qquad
 \bigl(\widehat{\pi_1(X)}\bigr)^F
     =\overline{\pi_1(L)}\cong\widehat{\mathbb Z}.
\end{equation}
This proves the actual critical-filling comparison in this region,
in both categories. It does not prove the global equality in B.106.

\begin{proof}
\emph{B.107.1. Separation of the actual punctures.}
Let $h_u=1/T_u$ and retain $d_u,\alpha_u$ from B.106.
Using $\overline u=u^{-1}$, direct expansion gives
\[
 |q-\alpha_u|^2-|d_u|^2
       =(|\alpha_u|^2-1)(1-|q|^2)>0
       \qquad(q\in\mathscr B_<).
\]
Thus $|h_u|<1$ and $|T_u|>1$, with $T_u=\infty$ allowed.
The two branch punctures $qt^2=1$ likewise have $|t|>1$.
All moving punctures therefore lie outside the unit disk.
The retained punctures $t=1,-1$ lie on its boundary.
The whole inside region is exactly $\mathscr B_<\times\{|t|<1\}$,
with equivariant radial contraction to the selected fixed source.

\emph{B.107.2. An actual retracted amalgam.}
Use the homotopy quotient of $\mathscr Q_<$ by the literal
$t\mapsto-t$ action. At the fixed base point its fundamental
group is $E=\pi_1(\mathscr Q_<)\rtimes\langle\tau\rangle$,
where $\tau$ represents that involution. Cut along
\[
 \mathscr B_<\times
       \bigl(\{|t|=1\}\setminus\{1,-1\}\bigr).
\]
This hypersurface is closed in the actual open surface. To specify
a collar, put
\[
 m(q)=\max\bigl(\sqrt{|q|},\max_{u\in U}|h_u(q)|\bigr)<1,
 \qquad \delta(q)=(1-m(q))/4.
\]
For $t=re^{i\theta}$, with $\theta\notin\pi\mathbb Z$, take
$|r-1|<\delta(q)|\sin\theta|$. Every moving puncture has modulus
at least $1/m(q)>1+\delta(q)$, so this collar avoids the deleted
divisors. The angular shrink keeps its two components separate at
the omitted $1,-1$. It is invariant under $t\mapsto-t$.
Enlarging the two sides across this collar gives an invariant
open cover, with the inside still equivariantly star-shaped to $0$.
Its intersection retracts
equivariantly to the two boundary arcs, which the involution
exchanges freely. The intersection homotopy quotient thus has
the homotopy type of $\mathscr B_<$.

Put $C=\pi_1(\mathscr B_<)\cong\mathbb Z$. The inside vertex
group is $C\times C_2$. Denote the outside vertex group by $R$;
it retains all the actual deleted graphs, branch curve and critical
fibers. The boundary section $t=i$, with one fixed choice of lift
in $EC_2$, gives the edge map $j:C\to R$. Projection to $q$
is a retraction $r:R\to C$, with $rj=1$. Thus this edge map
is injective. Its inside map is the ordinary $C$ factor.
Van Kampen, with a path at the fixed base value to the boundary,
gives the based isomorphism
\begin{equation}\label{eq:B107-amalgam}
 E=(C\times C_2)*_C R.
\end{equation}
The inside $C$ factor is exactly the selected source image.
The collar and projection justify the injections; no asphericity
or finite presentation of $R$ is assumed.

\emph{B.107.3. Centralizers in this amalgam, including completion.}
We prove the needed assertion for any discrete retract $C\le R$.
In \eqref{eq:B107-amalgam}, the ordinary reduced-word normal form
gives $C_E(\tau)=C\times C_2$: a reduced word leaving the first
vertex group cannot conjugate $\tau$ back into it, since $\tau$
does not belong to the amalgamated subgroup.

For full completion, regard $r$ as an idempotent endomorphism of
$R$ using $j$. Given any finite-index normal $N_0\lhd R$, put
$N=N_0\cap r^{-1}(N_0)$. Then $N$ is finite-index normal,
$N\le N_0$, and $r(N)\le N$. These $r$-stable subgroups form
a cofinal directed family. The induced retraction on the finite
group $Q_N=R/N$ has image $D_N$ and kernel $K_N$, so
$Q_N=K_N\rtimes D_N$. The corresponding quotient of $E$ is
\begin{equation}\label{eq:B107-heads}
 E_N=(D_N\times C_2)*_{D_N}Q_N
       =(C_2*K_N)\rtimes D_N,
\end{equation}
where $D_N$ fixes $\tau$. The subgroup $C_2*K_N$ has finite
index, so its full completion embeds as the open subgroup
$C_2\amalg K_N$ in $\widehat E_N$.
Herfort and Ribes' Theorem~B' gives
$C_{C_2\amalg K_N}(\tau)=C_2$; its full-category proof, using
open-subgroup Kurosh and Lemma~3, has been reopened at this use
on printed pp.~159--160~\citep{A_HR1985}.
If $K_N=1$ the same assertion is immediate. Consequently
\[
 C_{\widehat E_N}(\tau)=C_2\times D_N.
\]

Every finite quotient of $E$ factors through some $E_N$: start
with the kernel of its restriction to $R$ and refine to the
displayed $r$-stable subgroup. Its map on $C\times C_2$ then
factors through $D_N\times C_2$ as well. Hence the maps
$\widehat E\to\widehat E_N$ jointly separate elements.
For an element centralizing $\tau$, its images are uniquely
$\tau^\epsilon d_N$, with compatible coordinates. The groups
$D_N$ induce the full topology of $C$: for every finite-index
normal $L\lhd C$, the subgroup $N=r^{-1}(L)$ is $r$-stable
and gives $D_N=C/L$. Thus the compatible $d_N$ give one
$c\in\widehat C$. Subtracting $\tau^\epsilon c$ makes every
$\widehat E_N$ image trivial, and therefore makes the original
element trivial. We have proved
\[
 C_{\widehat E}(\tau)=\widehat C\times C_2.
\]
The displayed copy embeds by the explicit retractions. This
argument assumes neither residual finiteness of $R$ nor a
full-topology assertion for an unspecified infinite-index kernel.

\emph{B.107.4. Actual selected source and the remaining transfer.}
The finite-index extension $E=\pi_1(\mathscr Q_<)\rtimes C_2$
completes to the corresponding split extension, because its
kernel carries the full induced topology. Intersecting both
centralizer equalities with that kernel proves
\eqref{eq:B107-local-fixed} for $\mathscr Q_<$.
The source injection and its full topology also follow directly
from the projection to $\mathscr B_<$.

The actual cover $Z_<\to\mathscr Q_<$ has degree two. Over the
selected source, $\sqrt{1-q}$ is single-valued on the unit disk,
so the inverse image has two components. The component through
$P$ is
\[
 v=1+\sqrt{1-q},\qquad w=1-\sqrt{1-q},
 \qquad \sqrt{1-0}=1.
\]
It maps isomorphically to $\mathscr B_<$. Thus the selected
source in $\pi_1(\mathscr Q_<)$ lies in the index-two subgroup
from $Z_<$. Intersecting the already proved equalities with this
subgroup, and then with its full open completion, proves
\eqref{eq:B107-local-fixed} for $Z_<$. Only this direction of
cover transfer is used: the restricted source monodromy is trivial.

The reciprocal-conjugate symmetry
$(q,t)\mapsto(1/\overline q,1/\overline t)$ gives an outer-region
statement for $D_\infty$ in $\mathscr Q$, with a base point there.
Its lift concerns the $F\sigma$ source. It does not give the
selected $F$-source equality in the outer region of $Z$.
Moreover the fiber groups near the unit-circle arcs may acquire
kernels when the inner or outer critical fibers are inserted.
Their injections into the two regions have not been established.
Therefore \eqref{eq:B107-amalgam}, which cuts in the fiber
direction within the inner region, is not a global amalgam along
the $q$ boundary. The global equality in B.106 and the subsequent
main-B joint-source condition remain Open.
\end{proof}
